% GENERATED FILE. Edit canonical lessons, then run npm run generate:books.
% canonical-source-hash: fnv1a64:41bc76b67a12368c
% canonical-lessons: ES-C746-acariciar, ES-C746-engano, ES-C746-ligar, ES-C746-coquetear, ES-C746-amante

\chapter{Engaños y coqueteos}
\label{ch:es-enganos-y-coqueteos}
\hlchaptermodality{746}

\begin{chapteropening}
\textbf{By the end of this chapter:} \emph{I can talk about flirting, chatting someone up, caresses and deceit.}

Complete the last lesson of chapter 746: I can talk about flirting, chatting someone up, caresses and deceit.
\end{chapteropening}

\section[acariciar]{acariciar — to stroke, to caress}

\label{lesson:ES-C746-acariciar}



\begin{quote}
Before the new one: say the Spanish for to betray, then the Spanish for betrayal.
\end{quote}

\subsection*{You'll want to know: acariciar}
\textbf{acariciar} — \textquotedblleft{}to stroke, to caress\textquotedblright{}.

\begin{quote}
Le gusta acariciar al gato.
\end{quote}

\begin{quote}
\emph{She likes stroking the cat.}
\end{quote}

\emph{Una caricia} — a caress.

\begin{cousinweb}
\emph{acariciar} is built on \emph{caricia}, which was borrowed from Italian \emph{carezza} (EtymDB), from \emph{caro} \textquotedblleft{}dear\textquotedblright{}, Latin \emph{cārus}.
\end{cousinweb}

\begin{grammarlens}[title={What you've built}]
One more word for telling, explaining and describing at B1.
\end{grammarlens}

\subsection*{Your turn}
\begin{itemize}
\raggedright
  \item \emph{Say it:} \emph{acariciar}
  \item \emph{Say it:} \emph{acariciar}, once more
  \item \emph{Say it:} a sentence of your own with \emph{acariciar}
  \item \emph{Recall:} say the Spanish for to betray, then the Spanish for betrayal, then say \emph{acariciar} again
\end{itemize}

\begin{tcolorbox}[breakable,colback=teal!4,colframe=teal!35!black,title={Before you move on}]
What does acariciar mean? (\textquotedblleft{}To stroke, to caress\textquotedblright{}.) Say it once more.
\end{tcolorbox}
\section[el engaño]{el engaño — the deception, the trick}

\label{lesson:ES-C746-engano}



\begin{quote}
Before the new one: say the Spanish for betrayal, then the Spanish for to stroke.
\end{quote}

\subsection*{You'll want to know: el engaño}
\textbf{el engaño} — \textquotedblleft{}the deception, the trick\textquotedblright{}.

\begin{quote}
La oferta era un engaño.
\end{quote}

\begin{quote}
\emph{The offer was a trick.}
\end{quote}

\begin{cousinweb}
\emph{engaño} is built on \emph{engañar} \textquotedblleft{}to deceive\textquotedblright{}, which is usually traced to spoken Latin \emph{ingannāre} \textquotedblleft{}to mock\textquotedblright{}.
\end{cousinweb}

\begin{grammarlens}[title={What you've built}]
One more word for telling, explaining and describing at B1.
\end{grammarlens}

\subsection*{Your turn}
\begin{itemize}
\raggedright
  \item \emph{Say it:} \emph{el engaño}
  \item \emph{Say it:} \emph{el engaño}, once more
  \item \emph{Say it:} a sentence of your own with \emph{el engaño}
  \item \emph{Recall:} say the Spanish for betrayal, then the Spanish for to stroke, then say \emph{el engaño} again
\end{itemize}

\begin{tcolorbox}[breakable,colback=teal!4,colframe=teal!35!black,title={Before you move on}]
What does el engaño mean? (\textquotedblleft{}The deception, the trick\textquotedblright{}.) Say it once more.
\end{tcolorbox}
\section[ligar]{ligar — to chat up, to pick up; to tie}

\label{lesson:ES-C746-ligar}



\begin{quote}
Before the new one: say the Spanish for to stroke, then the Spanish for the deception.
\end{quote}

\subsection*{You'll want to know: ligar}
\textbf{ligar} — \textquotedblleft{}to chat up, to pick up; to tie\textquotedblright{}.

\begin{quote}
Fue a la discoteca a ligar.
\end{quote}

\begin{quote}
\emph{He went to the club to pick someone up.}
\end{quote}

Colloquial, mainly in Spain: \emph{ligar con alguien}. Also to tie, to bind.

\begin{cousinweb}
\emph{ligar} is usually traced to Latin \emph{ligāre} \textquotedblleft{}to tie, to bind\textquotedblright{} (Lewis \& Short); the slang sense of flirting grew from the idea of tying.
\end{cousinweb}

\begin{grammarlens}[title={What you've built}]
One more word for telling, explaining and describing at B1.
\end{grammarlens}

\subsection*{Your turn}
\begin{itemize}
\raggedright
  \item \emph{Say it:} \emph{ligar}
  \item \emph{Say it:} \emph{ligar}, once more
  \item \emph{Say it:} a sentence of your own with \emph{ligar}
  \item \emph{Recall:} say the Spanish for to stroke, then the Spanish for the deception, then say \emph{ligar} again
\end{itemize}

\begin{tcolorbox}[breakable,colback=teal!4,colframe=teal!35!black,title={Before you move on}]
What does ligar mean? (\textquotedblleft{}To chat up, to pick up; to tie\textquotedblright{}.) Say it once more.
\end{tcolorbox}
\section[coquetear]{coquetear — to flirt}

\label{lesson:ES-C746-coquetear}



\begin{quote}
Before the new one: say the Spanish for the deception, then the Spanish for to chat up.
\end{quote}

\subsection*{You'll want to know: coquetear}
\textbf{coquetear} — \textquotedblleft{}to flirt\textquotedblright{}.

\begin{quote}
Coquetea con todo el mundo.
\end{quote}

\begin{quote}
\emph{She flirts with everyone.}
\end{quote}

\emph{Coqueto} — flirtatious; also vain about one's looks.

\begin{cousinweb}
\emph{coquetear} is built on \emph{coqueta}, borrowed from French \emph{coquette} (EtymDB), from \emph{coq} \textquotedblleft{}cockerel\textquotedblright{}: showing off like a strutting cock.
\end{cousinweb}

\begin{grammarlens}[title={What you've built}]
One more word for telling, explaining and describing at B1.
\end{grammarlens}

\subsection*{Your turn}
\begin{itemize}
\raggedright
  \item \emph{Say it:} \emph{coquetear}
  \item \emph{Say it:} \emph{coquetear}, once more
  \item \emph{Say it:} a sentence of your own with \emph{coquetear}
  \item \emph{Recall:} say the Spanish for the deception, then the Spanish for to chat up, then say \emph{coquetear} again
\end{itemize}

\begin{tcolorbox}[breakable,colback=teal!4,colframe=teal!35!black,title={Before you move on}]
What does coquetear mean? (\textquotedblleft{}To flirt\textquotedblright{}.) Say it once more.
\end{tcolorbox}
\section[el amante]{el amante — the lover; lover of something}

\label{lesson:ES-C746-amante}



\begin{quote}
Before the new one: say the Spanish for to chat up, then the Spanish for to flirt.
\end{quote}

\subsection*{You'll want to know: el amante}
\textbf{el amante} — \textquotedblleft{}the lover; lover of something\textquotedblright{}.

\begin{quote}
Es un amante de la música.
\end{quote}

\begin{quote}
\emph{He is a music lover.}
\end{quote}

\emph{La amante} for a woman. \emph{Amante de} — keen on.

\begin{cousinweb}
\emph{amante} comes from Latin \emph{amāns}, \emph{amantem} \textquotedblleft{}loving\textquotedblright{} (EtymDB), the present participle of \emph{amāre}.
\end{cousinweb}

\begin{grammarlens}[title={What you've built}]
One more word for telling, explaining and describing at B1.
\end{grammarlens}

\subsection*{Your turn}
\begin{itemize}
\raggedright
  \item \emph{Say it:} \emph{el amante}
  \item \emph{Say it:} \emph{el amante}, once more
  \item \emph{Say it:} a sentence of your own with \emph{el amante}
  \item \emph{Recall:} say the Spanish for to chat up, then the Spanish for to flirt, then say \emph{el amante} again
\end{itemize}

\begin{tcolorbox}[breakable,colback=teal!4,colframe=teal!35!black,title={Before you move on}]
What does el amante mean? (\textquotedblleft{}The lover; lover of something\textquotedblright{}.) Say it once more.
\end{tcolorbox}
